\subsection{Accionamiento manual y uniones}
\label{app:prototipo_A8}
El radio efectivo entre ejes es 165\,mm, medido perpendicularmente al árbol.
La manivela conserva el tramo cilíndrico recto 15\,mm y su forma anterior;
su posición axial cambia y se incluye en A.7. Para cada estado:
\begin{equation}
 F_h=T/r_h,\qquad r_h=0,165\ {\rm m}.
\end{equation}
La fuerza ideal máxima es 16,848\,N.
\begin{center}\begin{tabular}{ll}\toprule
Eficiencia mecánica hipotética & $F_{h,\max}$ [N]\\\midrule
1,0 & 16,848\\
0,8 & 21,060\\
0,6 & 28,081\\
\bottomrule\end{tabular}\end{center}

Estas eficiencias no son mediciones. De forma general
$F_h=(T+T_{\text{pérdidas}})/r_h$. Un atasco no queda acotado por el torque
de operación:$T_{\rm mano}=r_hF_{\rm mano}$; no se inventa una carga máxima
de mano ni se declara resistencia al atasco.
El esfuerzo nominal del tramo recto es
\begin{equation}
 \sigma_b=32T/(\pi d_b^3)=8,390\ {\rm MPa},\quad d_b=15\ {\rm mm}.
\end{equation}
No es el máximo del agujero o codo. Se conservan los pendientes de
material y unión del brazo, sin investigar nuevas alternativas.

\subsubsection*{Matriz--árbol y árbol--buje}
La tuerca inferior apoya directamente en matriz; la arandela superior
permanece. Para presión uniforme, el radio medio de cada cara anular es
$r_m=\tfrac23(r_o^3-r_i^3)/(r_o^2-r_i^2)$: se obtienen
$r_{m,\inf}=8,1077$ mm y
$r_{m,\sup}=9,5225$ mm.
La condición global de transmisión por fricción, si ese es el mecanismo
efectivo, es
\begin{equation}
 T\leq\mu(N_{\inf}r_{m,\inf}+N_{\sup}r_{m,\sup}).
\end{equation}
Una condición conservadora con el menor radio sería
$\mu\sum N_i\geq342,88$ N.
No se asigna una precarga o fricción no documentada. Las tuercas de familia
adecuada permiten definir ese apriete, pero el modelo CAD no lo acredita.
El FEM adherente supone que la unión resiste el torque; no verifica por
sí solo esa condición. La unión manivela--árbol también requiere un
mecanismo real de apriete o transmisión positiva; el agujero 12,7\,mm no
basta para demostrarlo. El buje de adaptación resuelve el diámetro del
rodamiento principal, pero no se identifica un bloqueo axial positivo
del árbol al buje en el CAD.

\subsubsection*{Conclusión funcional de esta revisión}
El CAD resuelve la holgura exterior de los 608 y añade un apoyo principal
con adaptación de diámetro. La fuerza de compactación y el esfuerzo
manual permanecen prácticamente iguales; la flexión del árbol y la huella
de apoyo de matriz sí cambian y fueron recalculadas. La viabilidad bajo
las hipótesis de diseño se mantiene, condicionada a materializar la
retención axial y transmisión de torque indicadas y a las condiciones
resistentes ya pendientes, sin afirmar funcionamiento experimental.
